%% ============================================================================
%% APPENDIX: Extended Content for IOGraphNet Paper
%% ============================================================================

\appendix

\section{Extended Experimental Results}
\label{app:extended-results}

\subsection{Hyperparameter Search}

Table~\ref{tab:hyperparams} reports the hyperparameter search results.
We searched over learning rate, hidden dimensions, number of attention heads, and dropout rate using grid search with 5-fold cross-validation on the training set.
The best configuration uses learning rate 0.001, 64 hidden dimensions, 4 attention heads, and dropout 0.2.

\begin{table}[ht]
\caption{Hyperparameter search results (Micro-F1 on validation)}
\label{tab:hyperparams}
\small
\begin{tabular}{@{}lcccc@{}}
\toprule
LR & Hidden & Heads & Drop. & F1 \\
\midrule
0.01 & 64 & 4 & 0.2 & \placeholder{0.82} \\
0.001 & 32 & 4 & 0.2 & \placeholder{0.84} \\
\textbf{0.001} & \textbf{64} & \textbf{4} & \textbf{0.2} & \placeholder{\textbf{0.86}} \\
0.001 & 128 & 4 & 0.2 & \placeholder{0.85} \\
0.001 & 64 & 2 & 0.2 & \placeholder{0.84} \\
0.001 & 64 & 8 & 0.2 & \placeholder{0.85} \\
0.001 & 64 & 4 & 0.1 & \placeholder{0.84} \\
0.001 & 64 & 4 & 0.3 & \placeholder{0.85} \\
0.0001 & 64 & 4 & 0.2 & \placeholder{0.83} \\
\bottomrule
\end{tabular}
\end{table}

\subsection{Additional Ablation Studies}

\textbf{Number of GAT Layers.}
Table~\ref{tab:ablation-layers} shows the effect of varying the number of GAT layers.
Two layers provide the best trade-off between expressiveness and overfitting.
Deeper networks (3--4 layers) show slight degradation, likely due to over-smoothing where node representations become indistinguishable.

\begin{table}[ht]
\caption{Effect of number of GAT layers}
\label{tab:ablation-layers}
\small
\begin{tabular}{@{}lcc@{}}
\toprule
Layers & Micro-F1 & Training Time \\
\midrule
1 & \placeholder{0.82} & \placeholder{15} min \\
\textbf{2} & \placeholder{\textbf{0.86}} & \placeholder{20} min \\
3 & \placeholder{0.85} & \placeholder{28} min \\
4 & \placeholder{0.84} & \placeholder{35} min \\
\bottomrule
\end{tabular}
\end{table}

\textbf{Graph Pooling Methods.}
We compare mean pooling with alternative aggregation strategies (Table~\ref{tab:ablation-pooling}).
Mean pooling slightly outperforms max and sum pooling.
Attention-based pooling (Set2Set) provides marginal improvement but at higher computational cost.

\begin{table}[ht]
\caption{Effect of graph pooling method}
\label{tab:ablation-pooling}
\small
\begin{tabular}{@{}lcc@{}}
\toprule
Pooling & Micro-F1 & Inference Time \\
\midrule
\textbf{Mean} & \placeholder{\textbf{0.86}} & \placeholder{5} ms \\
Max & \placeholder{0.84} & \placeholder{5} ms \\
Sum & \placeholder{0.83} & \placeholder{5} ms \\
Set2Set & \placeholder{0.87} & \placeholder{12} ms \\
\bottomrule
\end{tabular}
\end{table}

\section{Dataset Details}
\label{app:dataset}

\subsection{Polaris System Characteristics}

ALCF Polaris is a 560-node HPE Apollo Gen10+ system with:
\begin{itemize}
    \item \textbf{Compute:} AMD EPYC Milan CPUs (32 cores), 4$\times$ NVIDIA A100 GPUs per node
    \item \textbf{Memory:} 512 GB DDR4 RAM, 160 GB HBM2e GPU memory per node
    \item \textbf{Storage:} Eagle Lustre file system (100 PB, 650 GB/s peak bandwidth)
    \item \textbf{Network:} HPE Slingshot interconnect (200 Gb/s per node)
\end{itemize}

\subsection{Polaris Log Statistics}

Table~\ref{tab:polaris-stats} provides detailed statistics for the Polaris production dataset.

\begin{table}[ht]
\caption{Polaris production dataset statistics}
\label{tab:polaris-stats}
\small
\begin{tabular}{@{}lr@{}}
\toprule
Statistic & Value \\
\midrule
Total logs & 1,378,117 \\
Date range & Apr 2024 -- Jan 2026 \\
Unique users & \placeholder{X} \\
Unique projects & \placeholder{Y} \\
\midrule
Logs with POSIX module & \placeholder{X\%} \\
Logs with MPI-IO module & \placeholder{Y\%} \\
Logs with STDIO module & \placeholder{Z\%} \\
Logs with HDF5 module & \placeholder{W\%} \\
\midrule
Median files per job & \placeholder{12} \\
Mean files per job & \placeholder{45} \\
Max files per job & \placeholder{10,000+} \\
\midrule
Median job duration & \placeholder{X} min \\
Median I/O bytes & \placeholder{Y} GB \\
\bottomrule
\end{tabular}
\end{table}

\subsection{Benchmark Configuration Details}

Table~\ref{tab:benchmark-configs} details the benchmark configurations used to generate ground-truth labels.

\begin{table*}[ht]
\caption{Benchmark configurations for ground-truth label generation (8 bottleneck categories)}
\label{tab:benchmark-configs}
\small
\begin{tabular}{@{}llll@{}}
\toprule
Bottleneck Type & Benchmark & Configuration & Label Criteria \\
\midrule
Metadata Storm & mdtest & 10K files/proc, stat-heavy & Meta time $>$ data time, opens/ops $>$10 \\
Small I/O & IOR & \texttt{-t 64 -b 64} (64-byte transfers) & $>$80\% ops $<$1KB \\
Data Imbalance & Custom MPI & Rank 0 writes 10$\times$ more; sleep injection & CV(bytes) $>$ 0.5, max/median time $>$ 2 \\
Write Saturation & IOR & Large sequential writes, high volume & High bytes, low bandwidth vs peak \\
Read Saturation & IOR & Large sequential reads, high volume & High bytes, low bandwidth vs peak \\
Interface Misuse & IOR & POSIX mode on shared file & Shared file, high indep ops, zero collective \\
File-per-Process & mdtest/Custom & 1000+ files created per rank & Files created $>$ 100 $\times$ ranks \\
Healthy & IOR & Optimized collective I/O & High bandwidth, balanced access \\
\bottomrule
\end{tabular}
\end{table*}

\subsection{Feature Normalization}

All node features are normalized using z-score standardization computed on the training set:
\begin{equation}
x'_i = \frac{x_i - \mu_i}{\sigma_i + \epsilon}
\end{equation}
where $\mu_i$ and $\sigma_i$ are the mean and standard deviation of feature $i$ across all training samples, and $\epsilon = 10^{-8}$ prevents division by zero.

For features with heavy-tailed distributions (e.g., byte counts, operation counts), we apply log transformation before normalization:
\begin{equation}
x'_i = \frac{\log(1 + x_i) - \mu_i}{\sigma_i + \epsilon}
\end{equation}

\section{Implementation Details}
\label{app:implementation}

\subsection{Structure-Aware Baseline Features}

To ensure fair comparison between IOGraphNet and flat-feature baselines, we augment XGBoost, Random Forest, and MLP with graph-derived statistical features.
Table~\ref{tab:baseline-features} lists these features, which allow baselines to access summarized structural information without the graph itself.

\begin{table}[ht]
\caption{Structure-aware features for baseline models}
\label{tab:baseline-features}
\small
\resizebox{\columnwidth}{!}{%
\begin{tabular}{@{}llp{4.5cm}@{}}
\toprule
\textbf{Category} & \textbf{Feature} & \textbf{Topological Intuition} \\
\midrule
\multirow{2}{*}{Global} & Total Ops (Read/Write/Meta) & Overall job intensity (AIIO-style) \\
 & Total Bytes (R/W) & Data volume \\
\midrule
\multirow{3}{*}{Graph (Node)} & Max/Mean Rank Degree & Fan-out: metadata stress \\
 & Max/Mean File Degree & Fan-in: contention/locking \\
 & Degree Std.\ Dev. & Load imbalance across nodes \\
\midrule
\multirow{3}{*}{Graph (Edge)} & Max Concurrency Weight & Peak contention on single file \\
 & Mean Concurrency Weight & General system pressure \\
 & Edge Weight Skewness & Stragglers vs.\ uniform I/O \\
\midrule
\multirow{2}{*}{Imbalance} & Gini Coefficient (bytes) & Data distribution inequality \\
 & Max/Mean Time Ratio & Straggler detection \\
\bottomrule
\end{tabular}%
}
\end{table}

\textbf{Key Insight:} If IOGraphNet outperforms structure-aware baselines, it proves that \emph{message-passing} (local neighborhood aggregation) provides value beyond statistical summaries.
The GNN can identify \emph{which specific node} is the outlier, while baselines only see ``average degree = 5.''

\subsection{Feature List}

Table~\ref{tab:features} lists all features extracted from Darshan logs for each file node.

\begin{table}[ht]
\caption{Node features extracted from Darshan logs (partial list)}
\label{tab:features}
\small
\begin{tabular}{@{}lll@{}}
\toprule
Category & Feature & Description \\
\midrule
\multirow{6}{*}{Operations} & read\_count & Number of read operations \\
& write\_count & Number of write operations \\
& bytes\_read & Total bytes read \\
& bytes\_written & Total bytes written \\
& seq\_reads & Sequential read count \\
& seq\_writes & Sequential write count \\
\midrule
\multirow{5}{*}{Access Size} & size\_0\_100 & Ops with size 0--100B \\
& size\_100\_1K & Ops with size 100B--1KB \\
& size\_1K\_10K & Ops with size 1KB--10KB \\
& ... & ... \\
& size\_1G\_plus & Ops with size $>$1GB \\
\midrule
\multirow{4}{*}{Timing} & read\_time & Total read time (s) \\
& write\_time & Total write time (s) \\
& meta\_time & Metadata operation time \\
& io\_time & Total I/O time \\
\midrule
\multirow{3}{*}{Derived} & avg\_read\_size & bytes\_read / read\_count \\
& bandwidth & bytes / time \\
& meta\_ratio & meta\_time / io\_time \\
\bottomrule
\end{tabular}
\end{table}

\subsection{Graph Construction Algorithm}

Algorithm~\ref{alg:graph-construction} describes the graph construction procedure.

\begin{algorithm}[ht]
\caption{Graph Construction from Darshan Log}
\label{alg:graph-construction}
\begin{algorithmic}[1]
\REQUIRE Darshan log $L$ with file records $\{r_1, ..., r_n\}$
\ENSURE Graph $G = (V, E, X)$
\STATE Initialize $V \leftarrow \emptyset$, $E \leftarrow \emptyset$
\FOR{each record $r_i \in L$}
    \STATE Add node $v_i$ to $V$
    \STATE Extract features $\mathbf{x}_i$ from $r_i$
\ENDFOR
\STATE // Add same-rank edges
\FOR{each pair $(v_i, v_j)$ where $i < j$}
    \IF{$\text{rank}(r_i) = \text{rank}(r_j)$}
        \STATE Add edge $(v_i, v_j)$ to $E$
    \ENDIF
\ENDFOR
\STATE // Add same-file edges
\FOR{each pair $(v_i, v_j)$ where $i < j$}
    \IF{$\text{filename}(r_i) = \text{filename}(r_j)$}
        \STATE Add edge $(v_i, v_j)$ to $E$
    \ENDIF
\ENDFOR
\STATE // Compute edge weights (concurrency)
\FOR{each edge $(v_i, v_j) \in E$}
    \STATE $w_{ij} \leftarrow \text{IoU}(\text{time\_window}(r_i), \text{time\_window}(r_j))$
\ENDFOR
\STATE // Fallback for sparse graphs
\IF{$|E| < 1$}
    \STATE Add edges between all node pairs
\ENDIF
\STATE Normalize features: $X \leftarrow \text{normalize}([\mathbf{x}_1, ..., \mathbf{x}_n])$
\RETURN $G = (V, E, X)$
\end{algorithmic}
\end{algorithm}

\subsection{Model Hyperparameters}

Table~\ref{tab:model-hyperparams} lists all model hyperparameters with justifications.

\begin{table}[ht]
\caption{IOGraphNet hyperparameters}
\label{tab:model-hyperparams}
\small
\begin{tabular}{@{}lll@{}}
\toprule
Parameter & Value & Justification \\
\midrule
GAT layers & 2 & Avoids over-smoothing \\
Hidden dim & 64 & Balance capacity/speed \\
Attention heads & 4 & Multi-perspective attention \\
Dropout & 0.2 & Regularization \\
Learning rate & 0.001 & Standard for Adam \\
Weight decay & $10^{-5}$ & L2 regularization \\
Batch size & 32 & GPU memory constraint \\
Max epochs & 100 & With early stopping \\
Early stop patience & 10 & Prevent overfitting \\
Activation & ELU & Smooth gradients \\
Pooling & Mean & Permutation invariant \\
\bottomrule
\end{tabular}
\end{table}

\section{Reproducibility}
\label{app:reproducibility}

\subsection{Code and Data Availability}

Upon acceptance, we will release:
\begin{itemize}
    \item \textbf{Code:} Complete implementation including data processing, model training, and evaluation scripts (GitHub repository, MIT license)
    \item \textbf{Benchmark dataset:} Darshan logs from IOR, mdtest, and DLIO with ground-truth labels
    \item \textbf{Trained models:} Pre-trained IOGraphNet weights for inference
    \item \textbf{Documentation:} Installation guide, usage examples, and API documentation
\end{itemize}

The Polaris production logs cannot be released due to ALCF data policies, but we provide detailed statistics (Appendix~\ref{app:dataset}) to enable replication on other systems.

\subsection{Hardware Specifications}

Experiments were conducted on:
\begin{itemize}
    \item \textbf{Training:} Single NVIDIA A100 (80 GB HBM2e)
    \item \textbf{CPU:} AMD EPYC 7763 (64 cores)
    \item \textbf{Memory:} 512 GB DDR4
    \item \textbf{Software:} Python 3.10, PyTorch 2.0, PyTorch Geometric 2.3
\end{itemize}

\subsection{Training Procedure}

\begin{enumerate}
    \item \textbf{Data preparation:} Parse Darshan logs, extract features, construct graphs
    \item \textbf{Split:} 70\% train, 15\% validation, 15\% test (stratified by label distribution)
    \item \textbf{Training:} Adam optimizer, batch size 32, max 100 epochs
    \item \textbf{Early stopping:} Monitor validation Micro-F1, patience 10 epochs
    \item \textbf{Model selection:} Checkpoint with best validation Micro-F1
    \item \textbf{Evaluation:} Final metrics on held-out test set
\end{enumerate}

Total training time: approximately \placeholder{2} hours on A100 GPU.

\subsection{Random Seeds}

For reproducibility, we fix random seeds:
\begin{itemize}
    \item Python random: 42
    \item NumPy: 42
    \item PyTorch: 42
    \item CUDA: deterministic mode enabled
\end{itemize}

We report mean and standard deviation over 5 runs with different seeds in Table~\ref{tab:main-results}.

%% ============================================================================
%% ARTIFACT DESCRIPTION / ARTIFACT EVALUATION APPENDIX
%% (Optional for HPDC 2026, up to 2 pages)
%% ============================================================================

\section{Artifact Description}
\label{app:artifact-description}

\subsection{Abstract}

IOGraphNet artifacts include the complete source code implementation, pre-trained model weights, labeled benchmark datasets, and reproduction scripts for all experimental results presented in this paper.
The artifacts enable independent verification of our main claims: (1) GNN-based bottleneck detection outperforms flat-feature baselines, (2) self-supervised pre-training improves generalization, and (3) attention weights provide interpretable root-cause localization.

\subsection{Artifact Availability}

\begin{itemize}
    \item \textbf{Repository:} {\footnotesize\url{https://anonymous.4open.science/r/iographnet-hpdc26}}
    \item \textbf{License:} MIT License (open source)
    \item \textbf{DOI:} To be assigned upon acceptance (Zenodo)
\end{itemize}

\subsection{Artifact Contents}

{\small
\begin{tabular}{@{}ll@{}}
\toprule
\textbf{Component} & \textbf{Description} \\
\midrule
\texttt{src/} & Source code (models, data loading) \\
\texttt{scripts/} & Reproduction scripts for experiments \\
\texttt{models/} & Pre-trained model weights \\
\texttt{data/sample/} & Sample benchmark data \\
\texttt{configs/} & Hyperparameter configurations \\
\texttt{README.md} & Documentation and usage \\
\texttt{requirements.txt} & Python dependencies \\
\bottomrule
\end{tabular}
}

\subsection{Hardware Dependencies}

\begin{itemize}
    \item \textbf{GPU:} NVIDIA GPU with CUDA support (tested on A100, V100)
    \item \textbf{Memory:} 32 GB RAM minimum (64 GB recommended)
    \item \textbf{Storage:} 10 GB for code, models, and sample data
\end{itemize}

\subsection{Software Dependencies}

\begin{itemize}
    \item Python 3.9 or higher
    \item PyTorch 2.0 or higher
    \item PyTorch Geometric 2.3 or higher
    \item CUDA 11.7 or higher (for GPU acceleration)
    \item See \texttt{requirements.txt} for complete list
\end{itemize}

\subsection{Datasets}

\begin{itemize}
    \item \textbf{Benchmark dataset:} 2,000+ labeled Darshan logs from IOR, mdtest, and DLIO benchmarks (included in artifacts)
    \item \textbf{Pre-training data:} 1.37M Polaris production logs (not redistributable due to ALCF policies; statistics provided in Appendix~\ref{app:dataset})
\end{itemize}

\section{Artifact Evaluation}
\label{app:artifact-evaluation}

\subsection{Setup Instructions}

\begin{verbatim}
# Clone repository
git clone [repository-url]
cd iographnet

# Create virtual environment
python -m venv venv
source venv/bin/activate

# Install dependencies
pip install -r requirements.txt

# Verify installation
python -c "import torch_geometric; print('OK')"
\end{verbatim}

\subsection{Reproducing Main Results (Table 3)}

\begin{verbatim}
# Run evaluation with pre-trained model
python scripts/evaluate.py \
    --model models/finetuned.pt \
    --data data/benchmark/test.pt

# Expected output:
# Micro-F1: 0.86 (+/- 0.02)
# Macro-F1: 0.78 (+/- 0.03)
# Subset Accuracy: 0.72 (+/- 0.02)
\end{verbatim}

\textbf{Time estimate:} 5 minutes (evaluation only)

\subsection{Reproducing Ablation Studies (Table 4)}

\begin{verbatim}
# Run ablation experiments
python scripts/run_ablations.py --output results/

# Expected output: ablation_results.csv
# matching Table 4 values within tolerance
\end{verbatim}

\textbf{Time estimate:} 30 minutes (GPU), 2 hours (CPU)

\subsection{Training from Scratch}

\begin{verbatim}
# Full training pipeline
python scripts/train.py \
    --config configs/default.yaml \
    --data data/benchmark/ \
    --output models/trained.pt

# Expected: Micro-F1 ~ 0.86 on test set
\end{verbatim}

\textbf{Time estimate:} 2 hours (A100 GPU), 8 hours (V100 GPU)

\subsection{Validation Criteria}

Results are considered successfully reproduced if:
\begin{itemize}
    \item Micro-F1 is within $\pm$0.02 of reported values
    \item Relative ranking of methods matches Table 3
    \item Ablation trends match Table 4 (graph edges improve performance)
\end{itemize}

Minor variations are expected due to hardware differences and random initialization.
